Gostaria de agradecer, primeiramente, a Deus, por me permitir chegar até este momento, guiando meus passos, fortalecendo-me diante das dificuldades e iluminando meu caminho nos períodos mais desafiadores desta jornada.

Aos meus pais, a quem dedico minha mais profunda admiração e gratidão, por todo o amor, sacrifício, apoio e ensinamentos que foram fundamentais na construção da pessoa que sou hoje. Nada do que conquistei seria possível sem o exemplo, os valores e o incentivo que sempre recebi de vocês. À minha irmã, agradeço por sua presença constante, pelo apoio emocional, pela amizade e por todas as conversas que me ajudaram a seguir em frente nos momentos em que mais precisei.

À minha orientadora, Prof.a Dr.a Edna Dias, expresso minha mais profunda gratidão. Sua orientação, dedicação e confiança foram fundamentais para que este trabalho se tornasse realidade. Em um momento de grande incerteza e desamparo durante o mestrado, a senhora assumiu minha orientação e me ofereceu não apenas direcionamento acadêmico, mas também incentivo e apoio para que eu pudesse seguir em frente. Tenho plena consciência de que, sem sua ajuda, esta trajetória poderia ter tomado um rumo muito diferente.

Agradeço por cada ensinamento, por cada correção cuidadosa e até mesmo pelos necessários “puxões de orelha”, que contribuíram de forma decisiva para meu crescimento como pesquisador e como pessoa. Sou igualmente grato por toda a disponibilidade, paciência e generosidade demonstradas ao longo deste percurso, muitas vezes indo além do que seria esperado de uma orientadora. Seu compromisso com a formação de seus alunos e sua genuína preocupação com seu desenvolvimento deixaram uma marca que levarei comigo muito além desta dissertação.

Por fim, agradeço pela paciência diante das minhas dificuldades e pelo apoio constante durante todo o processo de aprendizagem. Esta conquista também é fruto da confiança que a senhora depositou em mim.

Não poderia deixar de agradecer às integrantes da minha banca examinadora, Prof.a Dr.a Carla Taciana Lima e Profa. Dra. Edilene Lobo, pela disponibilidade e dedicação demonstradas ao longo deste trabalho, desde o período de qualificação até sua versão final. Mais do que avaliadoras, suas contribuições foram fundamentais para o amadurecimento desta pesquisa.

Agradeço pelos conselhos, questionamentos e correções, que me ajudaram não apenas a definir com maior clareza os aspectos mais relevantes do estudo, mas também a compreender melhor sua proposta, sua relevância e seu potencial impacto para a sociedade. As reflexões que compartilharam serviram como referência constante durante o desenvolvimento desta dissertação, orientando minhas decisões e fortalecendo minha confiança nos caminhos adotados ao longo da pesquisa.

Meus mais sinceros agradecimentos pela atenção, generosidade e valiosas contribuições, que certamente enriqueceram este trabalho.

Ao Programa de Pós-Graduação em Informática (PPGI), da Universidade de Brasília, agradeço por todo o suporte acadêmico e institucional oferecido ao longo desta jornada. Estendo minha gratidão aos coordenadores, professores e demais profissionais que contribuíram para minha formação e permanência no programa.

Em especial, agradeço à Profa. Dra. Claudia Nalon, cuja ajuda foi fundamental em um momento decisivo do meu percurso no mestrado. Sua orientação e apoio foram essenciais para que eu pudesse permanecer no programa e encontrar uma nova orientação, além da constante disponibilidade para esclarecer dúvidas e auxiliar sempre que necessário.

Gostaria também de registrar um agradecimento especial à Profa. Dra. Alba Cristina e ao Prof. Dr. Marcelo Marotta, responsáveis por algumas das disciplinas mais marcantes da minha formação. Além da excelência acadêmica, admiro profundamente a didática com que conduziram suas aulas, tornando o aprendizado uma experiência enriquecedora e inspiradora.

Por fim, agradeço ao PPGI pelo apoio acadêmico e financeiro disponibilizado ao longo desta trajetória, criando as condições necessárias para o desenvolvimento desta pesquisa e para minha formação como pesquisador.